%% EXHALE (EXHALE): comparison with reference escape codes.
%% Boundary conditions, initial conditions, and solver design, side by
%% side with the relevant reference atmospheric-escape codes:
%%   - Caldiroli et al. 2021 (A&A 655, A30): the original ATES code
%%   - Salz et al. 2016 (A&A 586, A75): PLUTO-CLOUDY interface (TPCI)
%%   - Kubyshkina et al. 2018 (A&A 619, A151): grid code, scheme from
%%     Erkaev et al. 2016 (MNRAS 460, 1300)
%%   - Murray-Clay et al. 2009 (ApJ 693, 23): steady-state BVP archetype
%% Source PDFs under ../../references/.
%% Author: Kwang-Il Seon
\documentclass[english,a4paper]{kiseon_memo}
\synctex=1
\usepackage{booktabs}
\usepackage{array}
\usepackage{amsmath}
\usepackage{babel}
\emergencystretch=3em

\newcommand{\Rp}{$R_{\rm pl}$}
\newcommand{\Mdot}{$\dot M$}

\title{Boundary conditions, initial conditions, and solver:\\
EXHALE / EXHALE vs.\ ATES, Salz+2016,\\
Kubyshkina+2018, and Murray-Clay+2009}
\author{Kwang-Il Seon}
\date{2026-06-13}

\begin{document}
\maketitle

\begin{quote}\small\itshape
A focused comparison of the boundary conditions (BC), initial conditions
(IC), and steady-state solver of EXHALE / EXHALE against the
relevant 1-D atmospheric-escape codes: its own parent code ATES
(Caldiroli et al.\ 2021), the PLUTO--CLOUDY interface (TPCI) of Salz et
al.\ (2016), the grid code of Kubyshkina et al.\ (2018) (numerical
scheme from Erkaev et al.\ 2016), and the steady-state boundary-value
archetype of Murray-Clay et al.\ (2009). The aim is to place EXHALE's
design choices in context --- what it inherited from ATES, what it
added, which choices are field-standard, and which capabilities the
other codes have that EXHALE lacks. A companion design sketch
(\texttt{auto\_ic\_design.md}) follows up on the main actionable
finding: a Kubyshkina-style automatic IC selector. Full BC/IC details
of EXHALE itself are in \texttt{EXHALE\_BC\_and\_IC.tex}.
\end{quote}

\tableofcontents

\section{The three codes at a glance}

\begin{table}[h]
\centering\small
\begin{tabular}{@{}p{0.20\textwidth}p{0.24\textwidth}p{0.24\textwidth}p{0.24\textwidth}@{}}
\toprule
 & \textbf{Salz+2016 (TPCI)} & \textbf{Kubyshkina+2018} & \textbf{EXHALE} \\
\midrule
Hydro scheme & WENO3 + HLLC + RK3 (Godunov) & McCormack predictor--corrector (finite diff., no Riemann solver) & PLM/WENO3 + HLLC + RK (Godunov) \\
Source term & operator-split forward Euler & in-scheme & semi-implicit energy + Brent $T$-solve \\
Steady state & time relaxation only & time relaxation & time relaxation \emph{+} JFNK Newton finish \\
Microphysics & CLOUDY full non-equilibrium photoionization (separate auto grid, 600--1200 cells) & H+H$_2$ network, Ly$\alpha$+H$_3^+$ cooling, $\eta=15\%$ fixed & ionization equilibrium (H/He+metals, MINPACK), CHIANTI metal cooling, He\,2$^3$S, H($n{=}2$)/Ly$\alpha$ RT \\
Composition & H/He (metals/molecules in CLOUDY only) & pure H; \textbf{H$_2$ molecular at base} & H/He \textbf{+ 10 metals} (self-consistent EOS) \\
Grid & stretched 500 cells, 1--12/15\,\Rp, base $\Delta r{=}2{\times}10^{-4}$\,\Rp & 1\,\Rp\ $\to$ Roche radius & 500 cells + 2 ghost, 1\,\Rp\ $\to r_{\max}$ \\
Potential & planet+star+centrifugal (2-body) + radiation pressure + conduction & Roche (Erkaev 2007) + conduction & Roche / spherical \\
\bottomrule
\end{tabular}
\caption{High-level comparison. EXHALE and Salz share the same
Godunov-type numerics (WENO3+HLLC+RK); Kubyshkina uses a lower-order,
more diffusive finite-difference scheme optimized for mass-producing a
$\sim$7000-model grid.}
\end{table}

The headline structural fact is that \textbf{EXHALE and Salz are
numerical near-twins} --- same reconstruction, same Riemann solver, same
time integrator --- so their differences are in the BC/IC details and in
the microphysics, not in the hydrodynamics. Kubyshkina is the outlier:
a deliberately simpler, faster scheme whose strengths lie in robustness
across escape regimes and in grid-production throughput. Murray-Clay is
a different beast entirely --- a steady-state boundary-value solver
rather than a time-relaxation code (\S\ref{sec:mc}).

\section{Heritage: what EXHALE inherited from ATES (Caldiroli+2021)}
\label{sec:heritage}

EXHALE / EXHALE is a fork of ATES, so its BC/IC/solver baseline is
literally Caldiroli et al.\ (2021). Reading that paper makes the
inheritance explicit and isolates what is genuinely new in EXHALE.

\textbf{Inherited verbatim (or nearly so):}
\begin{itemize}
\item \emph{Lower BC.} Density and temperature fixed at $R_{\rm pl}$,
  total $n=10^{14}\,\mathrm{cm^{-3}}$ (the same value Salz uses), $T$ at
  the planet's equilibrium temperature. \emph{The base velocity valve:}
  ``the velocity in the ghost cells at the lower boundary was obtained
  through a zeroth-order extrapolation from the first computational cell
  in the case of inflow characteristics, while it was set to zero
  otherwise'' --- i.e.\ EXHALE's $\max(v_1,0)$ valve is the original
  ATES valve. Caldiroli even states ``ATES employs the same boundary
  conditions as TPCI at the planet radius,'' so the ATES--Salz BC
  equivalence noted in \S\ref{sec:bc} is asserted in the ATES paper
  itself.
\item \emph{Upper BC.} Zero-gradient (PLM) / linear extrapolation
  (ESWENO3) at the Roche-lobe radius --- exactly EXHALE's outer BC.
\item \emph{IC.} ``A fully neutral, isothermal sphere of H and He
  $\ldots$ with a linear initial velocity profile $v(t{=}0)=(r-1)/2$''
  --- this \emph{is} EXHALE's cold-hydrostatic default IC.
\item \emph{Solver core.} Godunov finite-volume, SSPRK3 time
  integrator, HLLC default (+Roe, LLF), PLM and ESWENO3 reconstruction,
  $C_{\rm CFL}=0.6$, 500-cell mixed uniform--stretched grid, and the
  \emph{two-stage PLM\,$\to$\,ESWENO3 relaxation} with the final
  tolerance $\Delta\dot M/\dot M<10^{-3}$. Ionization solved in
  photoionization equilibrium via the MINPACK dogleg method; ion
  advection only as a post-process. All of this is EXHALE's backbone.
\end{itemize}

\textbf{Added by EXHALE} (none of these are in Caldiroli+2021): the
trace-metal network and its self-consistent EOS contribution
(\texttt{eos\_metals}); He\,2$^3$S, H($n{=}2$) and Ly$\alpha$ RT;
CHIANTI metal-line cooling; the \emph{JFNK Newton steady-state finish}
(ATES has \emph{no} implicit steady-state solver --- it is
time-relaxation only); the semi-implicit energy + Brent $T$-solve; the
\emph{smooth valve} (\texttt{Valve eps}) and the \emph{transonic} and
\emph{hot-Parker} IC families; the dynamic base-pressure correction
$\delta p_{\rm bc}$. In short, EXHALE keeps ATES's hydro core and BC/IC
skeleton and adds the metal physics, the excited-H/Ly$\alpha$ physics,
and a genuine steady-state solver on top. Notably, Caldiroli+2021 makes
\emph{no mention of base oscillations or ``breathing''}; that
phenomenology and its mitigations (smooth valve, PP-adv guard) are
EXHALE-era findings.

\section{Boundary conditions}
\label{sec:bc}

\subsection{Lower boundary}

\paragraph{Salz.} Two quantities fixed at the bottom --- \emph{density
and pressure} --- and \emph{none} at the top, which the paper notes is
``the correct choice for a subsonic inflow and a supersonic outflow.''
The base is anchored at $n\approx10^{14}\,\mathrm{cm^{-3}}$,
$p\approx14\,\mu$bar (set from $T_{\rm eq}$). Crucially, \emph{the
boundary-cell velocity is copied from the first interior cell each step
but only inflow (wind-feeding) is allowed; otherwise it is set to zero}
--- a choice the paper explicitly introduces ``to dampen oscillations
that appear with steep density gradients at the lower boundary.''

This is \textbf{exactly EXHALE's one-way velocity valve}
$v_{\rm gh}=\max(v_1,0)$, adopted independently for the same reason (the
breathing base, \texttt{EXHALE\_BC\_and\_IC.tex}\,\S2.4). EXHALE layers
three refinements on top that Salz does not have: (i) the dynamic
pressure correction $\delta p_{\rm bc}$ that pins the base to $T=T_0$
exactly through the EOS; (ii) the smooth valve (\texttt{Valve eps}) that
makes the kink differentiable for the Newton solver; (iii) the metal
mass/electron terms in the base density and pressure
(\texttt{eos\_metals}, 2026-06-13).

\paragraph{Kubyshkina (scheme from Erkaev+2016).} Lower boundary at the
photospheric radius \Rp\ with \emph{temperature fixed to $T_{\rm eq}$
and pressure set from Fossati+2017}, composition \emph{molecular H$_2$}.
The underlying Erkaev+2016 model is explicit: \emph{five} Dirichlet
conditions at the base ($\tilde T{=}1$, $\tilde\rho{=}1$, and the
hydrogen speciation fixed to pure H$_2$), with the base placed at the
\emph{homopause} (base of the thermosphere, $n_{\rm H_2}=5\times10^{12}\,
\mathrm{cm^{-3}}$), not the photosphere; the velocity is \emph{not}
fixed at the base. (Note the potential differs between the two papers:
Erkaev+2016 uses a \emph{point-mass} planetary potential, while
Kubyshkina+2018 adds the Roche term.)

Versus EXHALE: EXHALE fixes $(\rho,p)$ with $p$ dynamically corrected so
that $T=T_0$, i.e.\ \emph{effectively the same pair $(\rho,T)$}. The
substantive difference is the \textbf{H$_2$ molecular base}: it places
the lower boundary deeper, in a near-hydrostatic molecular layer, which
is what lets Kubyshkina follow boil-off (low-gravity, nearly hydrostatic
but rapidly escaping) atmospheres smoothly. EXHALE's atomic H/He base
does not represent that layer and is correspondingly weaker there ---
the same regime where EXHALE's breathing base is hardest to converge.

\subsection{The base \emph{radius} is the real lever, not the base
values (Murray-Clay+2009)}
\label{sec:rmin}

Murray-Clay et al.\ (2009) is the standard reference for \emph{how
much the lower BC matters}, and their answer is sharp: within physically
realistic ranges, the mass-loss rate is \textbf{insensitive} to the base
density, temperature, and ionization fraction (so long as
$\tau_{\rm base}\gg1$, $T_{\rm base}\ll10^4$\,K, $f_+\ll1$) --- they call
the solution ``quasi-unique.'' What \Mdot\ \emph{is} strongly sensitive
to is the base \emph{radius} $r_{\min}$: $\dot M\propto R_{\rm pl}^3$, so
a 20\% change in $r_{\min}$ moves \Mdot\ by up to a factor of two. Salz
independently finds the same insensitivity to the boundary
\emph{temperature} ($<5\%$ change in \Mdot\ for a 40\% temperature
change).

This reframes the whole lower-BC discussion and applies directly to
EXHALE: the precise values of $\rho_{\rm bc}$, $T_0$, and the new metal
EOS terms are \emph{not} the lever (consistent with the
\texttt{eos\_metals} mass effect being ${\sim}1\%$ in \Mdot), whereas the
choice of base \emph{radius} $R_0$ is. This echoes EXHALE's own
WASP-121b experience, where moving the base from the 4-mbar transit
radius ($R_0=1.766$) to the correct 1-$\mu$bar level ($R_0=2.058$) was
what brought Case A into agreement with Huang+2023 --- the same
``base-radius is the lever'' effect Murray-Clay quantify. Practical
takeaway: spend effort getting $R_0$ (and the pressure level it
corresponds to) right; do not over-tune $\rho_{\rm bc}$ or $T_0$.

\subsection{Upper boundary}

All three are essentially identical: \textbf{supersonic free outflow},
imposed as zero-gradient / linear extrapolation with no fixed condition
(Kubyshkina: ``radial derivatives become zero'' at the Roche radius;
Salz: ``none at the top''; EXHALE: copy / WENO3 linear extrapolation).
The same well-posedness caveat applies to all three when the flow is not
actually supersonic at the outer radius (deep Roche-lobe overflow);
EXHALE documents this explicitly (\texttt{EXHALE\_BC\_and\_IC.tex}\,\S3).

\section{Initial conditions}

\begin{table}[h]
\centering\small
\begin{tabular}{@{}p{0.13\textwidth}p{0.26\textwidth}p{0.26\textwidth}p{0.26\textwidth}@{}}
\toprule
 & Salz & Kubyshkina & EXHALE \\
\midrule
Method & single fixed IC (tuned to a medium-\Mdot\ atmosphere), same for all planets & \textbf{per-planet automatic selection of an initial profile} & cold hydrostatic (default) / transonic / hot-Parker \\
Stated property & ``final steady state does not depend on the IC'' & auto-selection accelerates a 7000-model grid & same \Mdot\ regardless; convergence is momentum-($du$)-gated \\
Transient & two shock waves cross in the first hours, then settle & --- & breathing-base limit cycle \\
\bottomrule
\end{tabular}
\caption{Initial-condition strategies.}
\end{table}

Salz demonstrates IC-independence of the steady state and gets away with
a single hand-tuned IC, at the cost of an expensive shock-laden
relaxation. \textbf{Kubyshkina's automatic per-planet IC selection is
the key engineering idea}: choosing an appropriate starting atmosphere
at each grid point is what makes mass-producing 7000 models tractable
and what lets the code launch winds across the full escape-parameter
range, including the boil-off cases that do not launch from a cold
static start. EXHALE already has the \emph{ingredients} for this --- the
transonic and hot-Parker IC families exist precisely for ``won't launch
from hydrostatic'' cases --- but the choice is currently \emph{manual}
(input keys), whereas Kubyshkina automates it. This is the main
actionable gap; see \texttt{auto\_ic\_design.md}.

EXHALE's IC behavior is, however, the best-documented of the three: the
benchmark
(\texttt{initial\_\allowbreak condition\_\allowbreak benchmark.pdf})
quantifies that a
naive full-Parker density seed diverges (self-amplifying breathing)
while the warm-seed converges to the same \Mdot\ with no speedup ---
because convergence is gated by the slow momentum relaxation, not the
thermal/ionization relaxation a warm seed shortcuts.

\section{Solver: differences, strengths, weaknesses}

\subsection{Versus Salz (the close comparison)}

\textbf{Shared:} WENO3/HLLC/RK family, the one-way valve, IC-independent
steady state.

\textbf{EXHALE advantages:}
\begin{itemize}
\item \emph{JFNK Newton steady-state finish.} Salz relaxes in time only,
  so its convergence criterion is $\Delta\dot M/\dot M<0.1$ (10\%) with
  residual oscillations handled by time-averaging. EXHALE converges to
  the fixed point directly ($\|R\|_\infty\to0$), a much stricter
  standard.
\item \emph{Metals self-consistent in the hydro EOS} (2026-06-13). Salz's
  published runs keep metals/molecules inside CLOUDY's separate grid and
  run H/He-only hydro.
\item \emph{Semi-implicit energy + Brent $T$-solve} avoids the hot-root
  trap of stiff metal cooling.
\end{itemize}

\textbf{Salz advantages:}
\begin{itemize}
\item \emph{CLOUDY's full non-equilibrium photoionization with
  self-consistent radiative transfer and molecular/metal cooling} ---
  fundamentally more accurate microphysics than EXHALE's ``ionization
  equilibrium + prescribed cooling fits,'' especially in the
  radiative-equilibrium base layers.
\item Very high base resolution ($\Delta r\sim2{\times}10^{-4}$\,\Rp)
  suppresses the pressure-gradient cancellation error that drives
  low-\Mdot\ oscillations.
\item Cost: TPCI is expensive and has no steady-state solver.
\end{itemize}

\subsection{Versus Kubyshkina}

\textbf{EXHALE advantages:}
\begin{itemize}
\item \emph{HLLC Riemann + WENO3} vs.\ the McCormack predictor--corrector
  $\Rightarrow$ much lower numerical diffusion. (Salz attributes the
  grid-dependent \Mdot\ of older Lax--Friedrich-type schemes, e.g.\ Tian
  et al.\ 2005, to exactly this diffusion; EXHALE/Salz are free of it,
  Kubyshkina's scheme is comparatively diffusive.)
\item Far richer microphysics (metals, He\,2$^3$S, Balmer/Ly$\alpha$ RT)
  vs.\ H/H$_2$ + Ly$\alpha$ + H$_3^+$ with a fixed $\eta=15\%$.
\end{itemize}

\textbf{Kubyshkina advantages:}
\begin{itemize}
\item \emph{Robust across the full boil-off / blow-off range.} The
  molecular-H$_2$ base plus per-planet automatic IC let it span escape
  parameters $\Lambda\sim20$--$80$ and produce 7000 models --- precisely
  the low-gravity / nearly-hydrostatic regime where EXHALE's breathing
  base is hardest to converge (e.g.\ HD\,189733\,b).
\item Simple and fast: ideal for the large model grids used in planetary
  evolution calculations. EXHALE is more accurate per case but heavier
  and breathing-limited.
\end{itemize}

The underlying scheme (Erkaev+2016) is a two-step MacCormack
predictor--corrector, \emph{second-order} and explicitly noted as
``suitable for smooth solutions without sharp discontinuities'' --- i.e.\
it lacks the shock-capturing of EXHALE's HLLC. It reaches steady state by
time relaxation (no artificial viscosity stated), and the paper validates
it against the analytic Parker (1958) wind and against Murray-Clay+2009
($3.33$ vs.\ $3.3\times10^{10}\,\mathrm{g\,s^{-1}}$). So the
EXHALE-vs-Kubyshkina solver gap is concrete: a Godunov/HLLC/WENO3 scheme
versus a smooth-flow second-order finite-difference scheme.

\subsection{A different solver philosophy: steady-state BVP
(Murray-Clay+2009)}
\label{sec:mc}

Murray-Clay+2009 is worth a dedicated note because it is the principal
\emph{methodological alternative} to everything above: instead of
relaxing a time-dependent code to a stationary state, it solves the
\emph{steady-state} equations directly as a two-point boundary-value
problem. The two points are the base and the \emph{sonic point}; six BCs
are imposed (two Parker critical-point conditions at the sonic point,
four at the base). The interior ($R_{\rm pl}\to r_s$) is solved by a
Newton relaxation (\texttt{solvde}) on the finite-differenced ODEs, and
the exterior ($r_s\to R_{\rm Roche}$) by Bulirsch--Stoer integration,
iterated for consistency.

Two lessons for EXHALE:
\begin{itemize}
\item \emph{The sonic point is the natural place to anchor a steady
  solver.} Murray-Clay stress that ``for every transonic wind there are
  an infinite number of breeze solutions,'' and codes that do not
  enforce the critical-point conditions ``found, not surprisingly,
  breezes instead.'' EXHALE's JFNK Newton finish is the time-domain
  analogue of their relaxation, but it does not explicitly impose the
  sonic-point conditions --- a possible robustness upgrade for the
  hard-to-launch cases, and an alternative to the transonic IC.
\item \emph{Continuation/homotopy is how they launch hard cases.} They
  build the final solution by ``solving successively more complicated
  problems'': start from an isothermal wind, then add photoionization,
  heating, cooling, and tidal gravity ``one by one $\ldots$ in a diluted
  form first and gradually strengthened to full amplitude.'' This is a
  proven recipe for reaching a transonic solution that EXHALE does not
  currently use; it is noted as an alternative IC/continuation strategy
  in \texttt{auto\_ic\_design.md}.
\end{itemize}

The trade-off: a steady-state BVP solver like Murray-Clay's is elegant
and fast for a single transonic wind, but it bakes in the transonic
topology (it cannot represent a genuinely time-dependent or
non-launching atmosphere) and does not naturally carry a large reaction
network. EXHALE's time-relaxation-plus-Newton approach is more general
at the cost of the breathing-base difficulty.

\subsection{Wind-AE (Broome et al.\ 2025): MC09's successor, verified
in-tree}
\label{sec:windae}

\emph{Added 2026-06-13.} The Murray-Clay code did not stay an archetype:
its direct successor, \textbf{Wind-AE} (Broome et al.\ 2025, ApJ 995,
198; BSD-3, PyPI \texttt{wind\_ae}), is available in this workspace
under \texttt{wind-ae-main/}. It keeps the MC09 BVP relaxation core
(Numerical Recipes \texttt{solvde} from the base to the sonic point,
Bulirsch--Stoer outward) and adds atomic metals, multifrequency XUV with
X-ray secondary/K-shell ionization, a toggleable tidal $+$ centrifugal
potential, and --- importantly --- a Python wrapper that implements
exactly the continuation/homotopy machinery of \S\ref{sec:mc}
(\texttt{ramp\_to}, \texttt{ramp\_grav}, \texttt{ramp\_Ftot},
\texttt{polish\_bcs}). Limitations relative to EXHALE: one ionization
stage per species (no Mg\,\textsc{ii}/Fe\,\textsc{ii} line work), no
Fe/Ca cooling, substellar-point streamline geometry, and integration
stops near the Coriolis turning radius. A Fortran port is planned in
\texttt{md\_backup/wind\_ae\_fortran\_port\_plan.md}; it is not
\emph{required} for use --- the C core compiles and runs on this Linux
machine as-is ($\sim$14\,s per solve) and the wrapper works under
Python 3.11 --- but remains under consideration for code
completeness/uniformity of the in-house tool set.

The following was verified hands-on (2026-06-13), ramping from the
bundled MC09 fiducial to EXHALE's own \texttt{input.inp} parameters with
matched integrated XUV flux (caveat: Wind-AE's bundled multifrequency
spectrum vs.\ EXHALE's power law; substellar vs.\ \texttt{Rate/2}
heating geometry — factor-of-a-few agreement is the realistic target):

\begin{itemize}
\item \textbf{HD\,209458b} (EXHALE values $0.720\,M_J$, $1.401\,R_J$,
  $F_{\rm XUV}=1081$\,erg\,cm$^{-2}$\,s$^{-1}$): converges in seconds;
  $T_{\rm max}=8\,640$\,K, substellar $4\pi r^2\rho v = 10^{10.78}$
  g\,s$^{-1}$ at the sonic point ($r_s=3.49\,R_{\rm pl}$). EXHALE's own
  metals-off run gives a flat $4\pi r^2\rho v = 10^{10.84}$ and
  $T_{\rm max}=8\,300$\,K --- 15\% / 4\% agreement from a completely
  independent method.
\item \textbf{HD\,189733b} (the planet where EXHALE's time marching
  limit-cycles at $du\sim2$ and never converges): the BVP relaxation
  \emph{converges cleanly} ($T_{\rm max}=11\,300$\,K,
  $r_s=3.83\,R_{\rm pl}$, global $\dot M = 10^{10.2}$ with their 0.33
  substellar-to-global factor). \textbf{A steady transonic solution
  exists for this planet.} This is the missing experiment from the
  2026-06-08 convergence study: the HD\,189733b oscillation is a
  property of EXHALE's time-marching $+$ base-BC numerics (acoustic
  reflection), not of the physics. The NSCBC / sonic-anchored-Newton
  routes are therefore worth their cost for this regime.
\item \textbf{WASP-121b} (near-RLOF, $F_{\rm XUV}\sim1.6\times10^{6}$):
  the \texttt{ramp\_to} continuation stalls around half the target flux
  (step size driven to $\sim$1\%); the deep-RLOF $+$ extreme-XUV corner
  appears not (easily) reachable --- consistent with the sonic-anchored
  formulation's difficulty when the critical point approaches L1.
  EXHALE's Roche machinery (Phase 4) has no replacement here.
\end{itemize}

\paragraph{Wind-AE as an EXHALE initial condition (tested).} A
$\sim$150-line converter (\texttt{src/utils/windae\_to\_exhale\_ic.py})
interpolates a \texttt{windsoln.csv} onto the EXHALE grid
(from \texttt{IC\_dump.txt}), blends hydrostatically down to the EXHALE
base anchor below Wind-AE's $R_{\rm min}\simeq1.06\,R_{\rm pl}$, extends
flux-conservingly above its $R_{\rm max}$, and writes EOS-consistent
\texttt{Hydro\_ioniz\_IC.txt} / \texttt{Ion\_species\_IC.txt} for
\texttt{Load IC}. Test on HD\,209458b (H/He, spherical, $10\,R_{\rm
pl}$, two-stage PLM$\to$WENO3, $du_{\rm th}=10^{-3}$):
\begin{itemize}
\item \emph{Hazard first:} with plain marching the run ``converges''
  \textbf{in 3 steps} --- a false stop. EXHALE's $du$ criterion measures
  the spatial spread of $\rho v r^2$, which a BVP solution satisfies
  \emph{by construction} ($du=2.5\times10^{-4}$ at load); the state is
  still Wind-AE's solution, not EXHALE's fixed point. Any Wind-AE-seeded
  run must use \texttt{Solver: Newton} (or an $R$-based gate; the same
  trap as the $\sim$5\% $du$-stop spread documented in
  \texttt{auto\_ic\_design.md} \S6.1, in its most extreme form).
\item \emph{Genuine value second:} under Newton-mode marching (no $du$
  stop) the loaded state evolves smoothly toward EXHALE's own attractor
  --- flux level $10^{10.78}\to10^{10.68}$, spread
  $2.5\times10^{-4}\to6.6\times10^{-2}$ (peak) $\to2.1\times10^{-2}$ at
  146k steps, monotone, no NaN. The \textbf{cold-start control crashed}
  (NaN at the breathing base, $r\simeq1.18$, step 172\,469, still at
  $du=3.7$) in the same configuration. For this hard, weakly-driven case
  the Wind-AE warm start survives where the cold start dies.
\end{itemize}
In short: Wind-AE is the working realization of both \S\ref{sec:mc}
lessons (sonic-point anchoring \emph{and} continuation), usable today as
an independent $\dot M$ oracle and as a warm-start IC generator for
EXHALE's hardest convergence regime --- with the explicit caveat that
its solution must never be accepted through the bare $du$ stop.

\section{What EXHALE should take away}

\begin{enumerate}
\item \textbf{The valve is standard, and EXHALE's version is the most
  complete.} Salz independently adopted the same one-way valve for the
  same anti-oscillation reason; EXHALE's $\delta p_{\rm bc}$, smooth
  valve, and metal EOS sit on top of it. The lower BC is not a weak
  point.
\item \textbf{The Newton finish is a genuine, unique strength} --- only
  EXHALE converges to the fixed point rather than time-averaging. The
  real limitation is not the solver but the breathing base in
  low-\Mdot\ / low-gravity cases.
\item \textbf{The most worthwhile idea to adopt is Kubyshkina's
  automatic IC selection} (and, relatedly, a deeper / molecular base for
  boil-off). EXHALE already has the transonic and hot-Parker IC families
  and computes the escape parameter $b_0\equiv\Lambda = GM_{\rm pl}\,
  m_{\rm H}/(k_BT_0R_{\rm pl})$; what is missing is an automatic rule
  that picks the IC family from $b_0$ and the sonic-point location. This
  directly targets EXHALE's weakest regime. See
  \texttt{auto\_ic\_design.md}.
\item \textbf{The microphysics ceiling is Salz/CLOUDY.} The known
  limitations of EXHALE's equilibrium + fitted cooling (the
  \texttt{use\_2lev\_cool} branch, fine-structure floors, etc.\ in
  \texttt{TO\_BE\_DONE.md}) are exactly what a non-equilibrium RT
  treatment like TPCI would resolve --- a long-term direction, not a
  near-term change.
\item \textbf{Tune the base \emph{radius}, not the base values}
  (Murray-Clay+2009; \S\ref{sec:rmin}). \Mdot\ is quasi-insensitive to
  $\rho_{\rm bc}$, $T_0$, and the new metal-EOS terms but
  $\propto R_{\rm pl}^3$ in the base radius. This is reassuring for the
  $\sim1\%$ \texttt{eos\_metals} change and is the lesson behind
  EXHALE's own $R_0=1.766\to2.058$ correction for WASP-121b. Effort on
  $R_0$/base-pressure calibration pays off; effort on $\rho_{\rm bc}$/$T_0$
  does not.
\item \textbf{Consider anchoring the steady solver at the sonic point}
  (Murray-Clay+2009; \S\ref{sec:mc}). EXHALE's Newton finish does not
  enforce the Parker critical-point conditions; doing so (or using
  continuation/homotopy to reach the transonic branch) is a candidate
  robustness upgrade for the non-launching cases, complementary to the
  automatic IC selector.
\item \textbf{Wind-AE makes items 2 and 6 actionable today}
  (\S\ref{sec:windae}). The MC09 successor runs in-tree: it proved a
  steady HD\,189733b solution exists (the limit cycle is numerics, not
  physics), cross-checks $\dot M$ at the 15\% level on HD\,209458b, and
  its solutions work as warm-start ICs that survive where the cold start
  NaN-crashes --- provided the run is finished with
  \texttt{Solver: Newton}, never the bare $du$ stop.
\end{enumerate}

\end{document}
